\documentclass[11pt,a4paper]{article}
\usepackage[utf8]{inputenc}\usepackage[T1]{fontenc}\usepackage{lmodern}
\usepackage{amsmath,amssymb,amsthm}\usepackage{geometry}\usepackage{xcolor}
\usepackage{hyperref}\usepackage{microtype}\usepackage{enumitem}\usepackage{booktabs}
\geometry{margin=2.5cm}
\hypersetup{colorlinks=true,linkcolor=blue!55!black,citecolor=blue!55!black,
            urlcolor=blue!55!black}
\theoremstyle{plain}
\newtheorem{theorem}{Theorem}[section]
\newtheorem{lemma}[theorem]{Lemma}
\newtheorem{proposition}[theorem]{Proposition}
\newtheorem{corollary}[theorem]{Corollary}
\theoremstyle{definition}
\newtheorem{remark}[theorem]{Remark}
\newcommand{\val}{\operatorname{v}_2}
\newcommand{\step}{\operatorname{step}}
\newcommand{\CC}{C_{\mathrm{Collatz}}}
\sloppy

\title{\bfseries Three Exact Results Arising from Collatz-Type Dynamics:\\
Valuation Equidistribution, Fermat Bases,\\ and a Transcendental Density Constant}
\author{Joel Cruz Cabrera\\
\small Santa Cruz de Tenerife, Canary Islands, Spain\\
\small\texttt{tf.joel.work@gmail.com}\quad ORCID: 0009-0005-4048-1237}
\date{July 2026}

\begin{document}\maketitle

\begin{abstract}
The study of the accelerated Collatz map $\step(x)=(3x+1)/2^{\val(3x+1)}$ throws off
questions that are not about the Collatz conjecture at all. We isolate three of them and
settle them completely.
\emph{(i)} The valuations produced by the accelerated step obey an exact $2$-adic law: for
every odd base $a$, $\val(an+1)\ge k \iff n\equiv-a^{-1}\pmod{2^{k}}$, so $\val(an+1)$ is
$\mathrm{Geom}(1/2)$ with mean $2$ and variance $2$, independently of $a$; and the two
valuations $\val(3n\pm1)$ of the two-sided step are anti-correlated with covariance exactly
$-1$.
\emph{(ii)} The map $x\mapsto ax+1$ ($a$ odd) admits closed-form polynomial orbits if and only
if $a-1$ is a power of $2$, i.e.\ $a=2^{e}+1$ is a Fermat-type base.
\emph{(iii)} The density constant $\CC=\frac12-\sum_{n\ge2}2^{-\lfloor n\log_{2}3\rfloor}
=0.2862745023241087\ldots$ has a Sturmian binary expansion and is therefore transcendental;
so is $C_{a}=1-\sum_{n\ge1}2^{-\lfloor n\log_{2}a\rfloor}$ whenever $\log_{2}a$ is irrational.
Result (i), and the orbit formula of (ii), are formalised in Lean~4 and machine-checked with
kernel-only axioms;
(iii) is verified by exact computation and follows from known transcendence criteria.
\textbf{None of the three depends on the Collatz conjecture, and none is claimed to advance
it.} We also record, with explicit witnesses, a limitative fact that is easy to get wrong
experimentally: for no modulus $M$ is $\val(3x+1)$ a function of $x\bmod M$, so no finite
automaton on residues can model the dynamics.
\end{abstract}

\noindent\textbf{Keywords:} Collatz map; $2$-adic valuation; equidistribution; Sturmian
words; transcendence; Beatty sequences; formalised mathematics; Lean~4.

\noindent\textbf{MSC 2020:} 11B83; 11A07; 11S05; 11J81; 68R15; 68V20; 37P05.

\section{Introduction}

For odd $x$ the \emph{accelerated} (Syracuse) Collatz step is
\begin{equation}\label{eq:step}
  \step(x)=\frac{3x+1}{2^{\val(3x+1)}},
\end{equation}
where $\val$ denotes the $2$-adic valuation. Whether every orbit reaches $1$ is the Collatz
conjecture, open since Collatz's 1937 formulation; see Lagarias~\cite{lagarias,lagarias-book}
for the standard survey, Terras~\cite{terras} for the classical density result, and
Tao~\cite{tao} for the strongest current ``almost all'' theorem. \textbf{This paper does not
address the conjecture.}

What it does address is a cluster of questions that arise while studying \eqref{eq:step} but
whose statements never mention it. Each is settled here, and each stands whether or not the
conjecture is true:

\begin{enumerate}[label=(\roman*),nosep]
  \item an exact equidistribution law for the valuations $\val(3^{r}m-1)$
        (Section~\ref{sec:geom});
  \item a characterisation of the bases $a$ for which $ax+1$ has closed-form orbits
        (Section~\ref{sec:fermat});
  \item the transcendence of a density constant attached to the map
        (Section~\ref{sec:trans}).
\end{enumerate}

\paragraph{Method, and why this is experimental mathematics.} All three began as
computational observations and were then either proved outright or reduced to a known
criterion. Two of the three are additionally \emph{formalised in Lean~4} and checked by its
kernel with no \texttt{sorry} and no custom axioms, which we regard as the appropriate
standard of evidence when the source of a conjecture is a machine. Section~\ref{sec:verif}
documents the verification, and Section~\ref{sec:nofinite} records --- with explicit
counterexamples --- a false reduction that is easy to reach experimentally and that we
believe deserves to be stated plainly in the literature.

\section{An exact equidistribution law}\label{sec:geom}

Throughout this section $m$ ranges over the odd residues modulo $2^{k}$, of which there are
$2^{k-1}$.

\begin{lemma}[Invertibility]\label{lem:inv}
$3$ is a unit in $\mathbb{Z}/2^{k}\mathbb{Z}$ for every $k\ge1$; hence $3^{r}$ is a unit and
$m\mapsto3^{r}m$ permutes $(\mathbb{Z}/2^{k}\mathbb{Z})^{\times}$.
\end{lemma}
\begin{proof}
$\gcd(3,2)=1$, so $3$ is invertible modulo every power of $2$. Units form a group, so $3^{r}$
is a unit and multiplication by it is a bijection of the unit group.
\end{proof}

\begin{theorem}[Exact equidistribution]\label{thm:geom}
Let $r,k\ge1$. There is \emph{exactly one} odd $m\in[0,2^{k})$ with $\val(3^{r}m-1)\ge k$.
\end{theorem}
\begin{proof}
$\val(3^{r}m-1)\ge k$ means $3^{r}m\equiv1\pmod{2^{k}}$, i.e.\ $m\equiv(3^{r})^{-1}
\pmod{2^{k}}$. By Lemma~\ref{lem:inv} this inverse exists, is unique modulo $2^{k}$, and is
odd (being a unit). Hence exactly one odd residue satisfies the condition.
\end{proof}

\begin{corollary}[Geometric law]\label{cor:geom}
Let $m$ be uniform on the odd residues modulo $2^{k}$. For $1\le j<k$,
\[
  \Pr\bigl[\val(3^{r}m-1)=j\bigr]=2^{-j},
\]
independently of $r$, and $\Pr[\val\ge k]=2^{-(k-1)}$. Letting $k\to\infty$, the valuation is
$\mathrm{Geom}(1/2)$ on $\{1,2,\dots\}$ with $\mathbb{E}[\val]=\sum_{j\ge1}j2^{-j}=2$.
\end{corollary}
\begin{proof}
By Theorem~\ref{thm:geom} the odd residues with $\val\ge j$ form a single class modulo
$2^{j}$, hence a fraction $2^{-(j-1)}$ of the odd residues; likewise $\Pr[\val\ge j+1]
=2^{-j}$. Subtract. Independence of $r$ is exactly Lemma~\ref{lem:inv}.
\end{proof}

\begin{remark}[Scope of the base]
The proof uses only that the base is odd (hence a $2$-adic unit) and that units form a group.
So for every odd $a$: for all $r,k\ge1$ exactly one odd $m\in[0,2^{k})$ has
$\val(a^{r}m-1)\ge k$, and the induced law is $\mathrm{Geom}(1/2)$ with mean $2$,
independently of $a$ and $r$.
\end{remark}

\begin{remark}[What it does \emph{not} give]\label{rem:scope}
Corollary~\ref{cor:geom} concerns \emph{uniformly distributed} residues. A single
deterministic orbit is not a random sample: its valuations are fixed, not drawn. The
corollary is therefore the exact form of the folklore heuristic ``$\mathbb{E}[\val]=2$, hence
expected factor $3/4<1$ per step'', but it implies nothing about any particular orbit. That a
negative average drift need not decide a deterministic orbit is shown concretely by the map
$n\mapsto\lfloor\sqrt2\,n+4\rfloor$ followed by halving, which Alth\"ofer and
Zipproth~\cite{az} prove diverges from every odd start.
\end{remark}

\subsection{The full statistics of the step valuations}\label{sub:stats}

The equidistribution law fixes the first moment. The same elementary machinery pins down the
whole joint law of the two valuations $v_{+}=\val(3n+1)$ and $v_{-}=\val(3n-1)$ produced by
one accelerated step, over $n$ uniform on the odd residues modulo $2^{K}$. Every statement
below is \emph{exact} (a rational, verified as such) and is machine-checked in the sense of
Section~\ref{sec:verif}; none uses floating point.

\begin{proposition}[Rotation law]\label{prop:rot}
For every odd $a$ and all $k\ge1$,
\[
  \val(an+1)\ge k \iff n\equiv -a^{-1}\pmod{2^{k}},
\]
a single residue class. Hence $\Pr[\val(an+1)=k]=2^{-k}$, independently of $a$: the base does
not change the law, it merely \emph{rotates} the class by $a^{-1}$.
\end{proposition}
\begin{proof}
$\val(an+1)\ge k$ means $an\equiv-1\pmod{2^{k}}$, i.e.\ $n\equiv-a^{-1}$, which exists and is
unique because $a$ is a $2$-adic unit. One class out of $2^{k-1}$ odd residues gives the
geometric weight; subtracting consecutive levels gives $2^{-k}$.
\end{proof}

\begin{proposition}[Second moment]\label{prop:var}
$\mathbb{E}[v_{+}^{2}]=6-(4K+6)2^{-K}$, so $\mathbb{E}[v_{+}^{2}]\to6$ and
$\operatorname{Var}(v_{+})\to 6-2^{2}=2$ --- exactly the mean and variance of $\mathrm{Geom}(1/2)$.
\end{proposition}

\begin{proposition}[The two valuations are anti-correlated]\label{prop:cov}
$\val(3n+1)$ and $\val(3n-1)$ satisfy $\min(v_{+},v_{-})=1$ identically: $3n+1$ and $3n-1$ are
consecutive even numbers, so exactly one is $\equiv2\pmod4$. Consequently
$v_{+}v_{-}=\max(v_{+},v_{-})$, whence $\mathbb{E}[v_{+}v_{-}]\to3$ and
\[
  \operatorname{Cov}(v_{+},v_{-})=\mathbb{E}[v_{+}v_{-}]-\mathbb{E}[v_{+}]\,\mathbb{E}[v_{-}]
  = 3-2\cdot2 = -1 .
\]
\end{proposition}
\begin{proof}
If $v_{+},v_{-}$ were independent the product mean would be $4$. Because one factor is always
$1$, the product equals the other factor, i.e.\ $\max(v_{+},v_{-})$; and
$\mathbb{E}[\max]=\mathbb{E}[v_{+}+v_{-}]-\mathbb{E}[\min]=4-1=3$. Subtracting the product of
the means gives $-1$.
\end{proof}

\noindent The anti-correlation is the quantitative face of the up/down dichotomy: whenever one
child of the two-sided step $n\mapsto\operatorname{odd\_part}(3n\pm1)$ rises sharply, the other
is \emph{forced} to valuation $1$. The dichotomy is usually stated qualitatively; here it has a
number, $\operatorname{Cov}=-1$.

\section{Which bases admit closed-form orbits}\label{sec:fermat}

Fix an odd $a\ge3$ and consider $\step_{a}(x)=\operatorname{odd\_part}(ax+1)$ on odd integers.
Every odd $x$ has a unique \emph{block} form
\begin{equation}\label{eq:block}
  x=2^{r+1}m-1,\qquad m\ \text{odd},\ r\ge0,
\end{equation}
in which $r=\val(x+1)-1$.

\begin{theorem}[One-step law and the characterisation]\label{thm:fermat}
Let $a\ge3$ be odd and write $a-1=2^{e}q$ with $q$ odd. For every odd $m$ and every $r\ge e$,
put $x=2^{\,r+1}m-1$. Then
\begin{equation}\label{eq:onestep}
  \val(ax+1)=e\qquad\text{and}\qquad \step_{a}(x)=a\,2^{\,r+1-e}m-q .
\end{equation}
Consequently the closed-form law $\step_{a}(x)=a\,2^{\,r+1-e}m-1$ holds \emph{if and only if}
$q=1$, i.e.\ iff $a-1$ is a power of $2$, i.e.\ iff $a=2^{e}+1$. In that case iteration gives
\begin{equation}\label{eq:orbit}
  x_{i}=a^{i}\,2^{\,r+1-ie}m-1\qquad\text{for every }i\text{ with }ie\le r .
\end{equation}
\end{theorem}

\begin{proof}
Compute directly:
\[
  ax+1=a\bigl(2^{\,r+1}m-1\bigr)+1=a\,2^{\,r+1}m-(a-1)
      =2^{e}\bigl(a\,2^{\,r+1-e}m-q\bigr).
\]
Since $r\ge e$ the bracket is (even)$-$(odd), hence odd; so $\val(ax+1)=e$ \emph{exactly} and
$\step_{a}(x)$ equals the bracket, which is \eqref{eq:onestep}. The image equals
$a2^{\,r+1-e}m-1$ precisely when $q=1$, which is both directions of the characterisation.

For \eqref{eq:orbit}, let $q=1$. The image $a2^{\,r+1-e}m-1$ again has the shape
\eqref{eq:block}, with odd parameter $m'=am$ and exponent $r'+1=r+1-e$; so the family is
preserved and induction on $i$ gives \eqref{eq:orbit} as long as $ie\le r$.
\end{proof}

\noindent Note that \eqref{eq:onestep} holds for \emph{every} odd $a$ --- it is the general
one-step law --- and the characterisation is the observation that its right-hand side lands
back in the family exactly when the odd part $q$ of $a-1$ is trivial.

\begin{remark}[Terminology]
The bases $a=2^{e}+1$ are $3,5,9,17,33,65,129,257,513,\dots$ We call them
\emph{Fermat-type} because the Fermat numbers $2^{2^{k}}+1$ are the sub-family $e=2^{k}$;
the two sets are \emph{not} equal --- $9$, $33$, $65$, $129$ and $513$ are of the form
$2^{e}+1$ without being Fermat numbers, and only $3,5,17,257,65537$ are Fermat primes.
\end{remark}

\begin{table}[ht]\centering
\begin{tabular}{c c c c}
\toprule
$a$ & $a-1=2^{e}q$ & $q$ & closed form?\\\midrule
$3$  & $2^{1}\cdot1$  & $1$ & yes (Collatz)\\
$5$  & $2^{2}\cdot1$  & $1$ & yes\\
$7$  & $2^{1}\cdot3$  & $3$ & no\\
$9$  & $2^{3}\cdot1$  & $1$ & yes\\
$11$ & $2^{1}\cdot5$  & $5$ & no\\
$13$ & $2^{2}\cdot3$  & $3$ & no\\
$15$ & $2^{1}\cdot7$  & $7$ & no\\
$17$ & $2^{4}\cdot1$  & $1$ & yes\\
$33$ & $2^{5}\cdot1$  & $1$ & yes\\
\bottomrule
\end{tabular}
\caption{The closed form holds exactly when the odd part $q$ of $a-1$ is $1$
(Theorem~\ref{thm:fermat}). The one-step law \eqref{eq:onestep} and the closed form were
checked independently of each other --- so the ``no'' column is a genuine test, not a
restatement of the hypothesis --- for all odd $a\le999$, $m\le200$ odd and $8$ consecutive
$r\ge e$. The bases with $q=1$ up to $999$ are $3,5,9,17,33,65,129,257,513$.}
\label{tab:fermat}
\end{table}

\begin{corollary}
Collatz ($a=3$) is the smallest Fermat-type base; the next are $5,9,17,33,\dots$ (all such
bases are odd, $2^{e}+1$ being odd for $e\ge1$). The
heuristic drift of $\step_{a}$ is $\lambda_{a}=\log a-\mathbb{E}[\val]\log2=\log a-2\log2$
(by Corollary~\ref{cor:geom}), negative iff $a<4$. Hence $a=3$ is the unique odd base on the
contracting side, and every odd $a\ge5$ is super-critical.
\end{corollary}
\noindent In view of Remark~\ref{rem:scope}, the corollary is a statement about averages and
decides no individual orbit; in particular it says nothing about whether $9x+1$ diverges from
$x=1$.

\section{A transcendental density constant}\label{sec:trans}

Counting the odd integers whose first block does not immediately shrink leads to
\begin{equation}\label{eq:def}
  \CC=\frac12-\sum_{n\ge2}2^{-\lfloor n\log_{2}3\rfloor}.
\end{equation}
Summing \eqref{eq:def} to $n=300$ in $80$-digit arithmetic (the tail is below $2^{-475}$):
\[
  \CC=0.2862745023241087\,4196\ldots
\]

\begin{lemma}[Beatty bookkeeping]\label{lem:beatty}
For every $n\in\mathbb{N}$, $\;n+\lfloor n\log_{2}(3/2)\rfloor=\lfloor n\log_{2}3\rfloor$.
\end{lemma}
\begin{proof}
$\log_{2}(3/2)=\log_{2}3-1$ and $\lfloor y-k\rfloor=\lfloor y\rfloor-k$ for $k\in\mathbb{Z}$.
\end{proof}
\noindent The identity is elementary; we record it only because it identifies the two Beatty
sequences that occur, and because the equivalent closed form
$1-\sum_{n\ge1}2^{-(n+\lfloor n\log_{2}(3/2)\rfloor)}=\CC$ is the shape that was formalised.

Set $\beta=\log_{2}3\in(1,2)$, irrational, with Beatty sequence
$B=(\lfloor n\beta\rfloor)_{n\ge1}$, and gap word $w=(w_{n})_{n\ge1}$,
$w_{n}=\lfloor n\beta\rfloor-\lfloor(n-1)\beta\rfloor$.

\begin{lemma}\label{lem:gap}
$w_{n}\in\{1,2\}$ for all $n$.
\end{lemma}
\begin{proof}
$\lfloor y+\beta\rfloor-\lfloor y\rfloor\in\{\lfloor\beta\rfloor,\lfloor\beta\rfloor+1\}
=\{1,2\}$ since $1<\beta<2$.
\end{proof}

\begin{lemma}[Sturmian]\label{lem:sturm}
$w$ is Sturmian: its factor complexity is $p(k)=k+1$ for every $k\ge1$.
\end{lemma}
\begin{proof}
By Lemma~\ref{lem:gap}, $w$ is the difference word of a Beatty sequence of irrational slope
$\beta$; such words are precisely the mechanical words of slope $\beta-\lfloor\beta\rfloor$,
which are Sturmian exactly because the slope is irrational
(Morse--Hedlund~\cite{mh}; see Lothaire~\cite{lothaire}, Ch.~2).
\end{proof}
\noindent Computing the factors of $w$ directly for $n\le4000$ returns alphabet $\{1,2\}$ and
$p(k)=k+1$ for all $k\le8$, as predicted.

Since the exponents in \eqref{eq:def} increase by $w_{n}\in\{1,2\}$, the binary digits of the
series are supported exactly on the Beatty set $\{\lfloor n\beta\rfloor\}$; the digit word of
$\CC$ is the characteristic word of that set, Sturmian by Lemma~\ref{lem:sturm}.

\begin{theorem}\label{thm:trans}
$\CC$ is transcendental.
\end{theorem}
\begin{proof}
The binary expansion of $\CC$ is Sturmian, hence aperiodic (so $\CC$ is irrational) with
factor complexity $p(k)=k+1$. A real number whose base-$b$ expansion is Sturmian is
transcendental by Ferenczi--Mauduit~\cite{fm}; alternatively, the complexity criterion of
Adamczewski--Bugeaud~\cite{ab} shows that an algebraic irrational cannot have an expansion of
complexity $p(k)=k+1$.
\end{proof}

\begin{corollary}
For every real $a>1$ with $\log_{2}a\notin\mathbb{Q}$, the constant
$C_{a}=1-\sum_{n\ge1}2^{-\lfloor n\log_{2}a\rfloor}$ is transcendental; in particular for
every integer $a$ that is not a power of $2$.
\end{corollary}
\begin{proof}
Identical, with $\beta=\log_{2}a$: only irrationality of $\beta$ is used, the gap alphabet
being $\{\lfloor\beta\rfloor,\lfloor\beta\rfloor+1\}$.
\end{proof}

\begin{remark}
The continued fraction of $\CC$ begins
\[
  [0;\,3,\,2,\,36,\,2056,\,70368878395648,\,\dots],
\]
the large early partial quotients reflecting exceptionally good dyadic-rational
approximations; only terms stable across two independent precisions are reported. We checked
$\CC$ against the CODATA physical constants and the standard mathematical constants and found
no match within $10^{-3}$; we know of no closed form for it, and it is transcendental by Theorem~\ref{thm:trans}. It is
not, however, a constant of a new kind: it belongs to the classical family
$\sum_{n\ge1} b^{-\lfloor n\alpha\rfloor}$, studied since B\"ohmer~\cite{bohmer}, who proved
transcendence when $\alpha$ has unbounded partial quotients, and later by
Danilov~\cite{danilov}, Davison~\cite{davison} and Adams--Davison~\cite{adamsdavison}, the last
of whom give the continued fraction of $(b-1)\sum_{n\ge1} b^{-\lfloor n\alpha\rfloor}$
explicitly for every irrational $\alpha$. What is specific here is only that this member of
the family, with $b=2$ and $\alpha=\log_2 3$, is the one the Collatz observable produces.
\end{remark}

\section{No finite-residue reduction}\label{sec:nofinite}

Experimentation with \eqref{eq:step} makes it tempting to conclude that the step is
determined by a residue of $x$ --- for instance modulo $24$, where the residue \emph{class}
does evolve deterministically. It is not, and we state this explicitly because the error
recurs.

\begin{proposition}\label{prop:nofinite}
For every modulus $M$, the valuation $\val(3x+1)$ is \emph{not} a function of $x\bmod M$.
\end{proposition}
\begin{proof}
Write $M=2^{s}M'$ with $M'$ odd, and let $c$ be the unique solution of $3c\equiv-1
\pmod{2^{s+1}}$; it is odd, since $c=-3^{-1}$ and $3^{-1}$ is a unit. Every $x\equiv c
\pmod{2^{s}}$ has $\val(3x+1)\ge s$, but the exact value is decided by the binary digits of
$x$ above position $s$, which the residue $x\bmod M$ does not see. Concretely, by the Chinese
Remainder Theorem choose odd $x,y$ with $x\equiv y\equiv c\pmod{2^{s}}$ and $x\equiv y
\pmod{M'}$ --- so $x\equiv y\pmod M$ --- while additionally $3x+1\equiv2^{s}\pmod{2^{s+1}}$
and $3y+1\equiv0\pmod{2^{s+1}}$; the two congruence systems are compatible because they
constrain $x$ and $y$ modulo $2^{s+1}$ and $M'$ at coprime places. Then $\val(3x+1)=s$ while
$\val(3y+1)\ge s+1$.
\end{proof}

\noindent The point is that $\val(3x+1)\ge s$ is decided by $x\bmod 2^{s}$, but $\val(3x+1)=s$
exactly is not: the valuation sequence needs unbounded $2$-adic precision, so the dynamics is
\emph{not} a finite automaton on residues.

\paragraph{The case $M=24$, explicitly.} Among the $2000$ odd $x<4000$:
\begin{itemize}[nosep]
  \item $500$ ($25\%$) have $\val(3x+1)\ne\val(3(x+24)+1)$; the smallest is $x=5$, $y=29$,
        with $\val(16)=4\ne3=\val(88)$.
  \item $1959$ ($98\%$) have $\step(x)\not\equiv\step(x+24)\pmod{24}$ --- so even the residue
        class of the image is not determined. The smallest witness is $x=1$, $y=25$:
        $\step(1)=1$ and $\step(25)=19$, and $1\not\equiv19\pmod{24}$.
\end{itemize}
An argument reducing the dynamics to finitely many residue classes therefore cannot be
correct. We stress a related trap: verifying a claimed identity on the $M$ residue
\emph{representatives} $x=1,3,\dots,M-1$ establishes nothing about the other members of each
class --- $x=1$ passes while $x=25$, in the same class, fails.

\section{Verification and formalisation}\label{sec:verif}

\begin{table}[ht]\centering
\small
\begin{tabular}{@{}l l l@{}}
\toprule
Result & Status & Evidence\\\midrule
Thm~\ref{thm:geom} & proved; machine-checked &
  Lean~4: \texttt{geometric\_distribution\_informal}\\
Cor.~\ref{cor:geom} & proved (counting) & from Thm~\ref{thm:geom}; recomputed exactly\\
Prop.~\ref{prop:rot} & proved; machine-checked & Lean~4: rotation law, $6$ bases certified\\
Prop.~\ref{prop:var} & proved & $\mathbb{E}[v_{+}^{2}]=6-(4K+6)2^{-K}$ exact, $K\le18$\\
Prop.~\ref{prop:cov} & proved & $\operatorname{Cov}=-1$ exact; $v_{+}v_{-}=\max$ ($0$ exc.)\\
Thm~\ref{thm:fermat} & proved; machine-checked (both ways) &
  Lean~4: \texttt{closed\_form\_iff\_fermat}\\
Lem.~\ref{lem:beatty} & proved (elementary) & one line; also in Lean\\
Lem.~\ref{lem:sturm} & proved (Morse--Hedlund) & $p(k)=k+1$ recomputed for $k\le8$\\
Thm~\ref{thm:trans} & proved via \cite{fm,ab} & $\CC$ to $80$ digits; expansion Sturmian\\
Prop.~\ref{prop:nofinite} & proved & $1959/2000$ witnesses at $M=24$, $x<4000$\\
\bottomrule
\end{tabular}
\caption{Status of each result. ``Machine-checked'' means Lean's kernel accepts the proof
with no \texttt{sorry} and only the ambient axioms \texttt{propext},
\texttt{Classical.choice}, \texttt{Quot.sound}.}
\end{table}

The Lean development builds with \texttt{lake build CollatzProof.KodamaCollatz
CollatzProof.KodamaFermat}, and \texttt{\#print axioms} is inspected for every theorem cited
above; all report exactly \texttt{[propext, Classical.choice, Quot.sound]}. The audit is kept
as \texttt{CollatzProof/AxiomAudit.lean} so that it re-runs on every build. Both halves of
Theorem~\ref{thm:fermat} are in the kernel: \texttt{step\_a\_general} is the one-step law
\eqref{eq:onestep} for every odd base, and \texttt{closed\_form\_iff\_fermat} is the
biconditional itself. The numerical statements (the value of $\CC$, the gap alphabet, the
factor complexity, Table~\ref{tab:fermat}, the witness counts) are regenerated by a single
script using exact integer and high-precision decimal arithmetic --- no floating point is used
where exactness is claimed.

\section{What is not claimed}

\begin{itemize}[nosep]
  \item No progress on the Collatz conjecture. Sections~\ref{sec:geom}--\ref{sec:trans} are
        independent of it, and Section~\ref{sec:nofinite} is limitative.
  \item No claim of novelty for Lemma~\ref{lem:beatty}, which is one line.
  \item Corollary~\ref{cor:geom} is about uniform residues, not about individual orbits
        (Remark~\ref{rem:scope}).
\end{itemize}

\section*{Acknowledgements}
The author thanks Ingo Alth\"ofer for maintaining an open problem list that prompted several
of these questions, and the Lean~4 and \emph{mathlib} communities. This work is AI-assisted;
every statement reported here was verified by exact computation or by Lean's kernel before
inclusion, and every candidate that failed verification was discarded.

\begin{thebibliography}{99}
\bibitem{ab} B.~Adamczewski and Y.~Bugeaud,
\emph{On the complexity of algebraic numbers I. Expansions in integer bases},
Ann. of Math. (2) \textbf{165} (2007), no.~2, 547--565.

\bibitem{az} I.~Alth\"ofer and T.~Zipproth,
\emph{Divergence in a Collatz-type recursion}, preprint, August 2025, \url{https://althofer.de/althoefer-zipproth-august-2025.pdf}.

\bibitem{fm} S.~Ferenczi and C.~Mauduit,
\emph{Transcendence of numbers with a low complexity expansion},
J. Number Theory \textbf{67} (1997), no.~2, 146--161.

\bibitem{lagarias} J.~C. Lagarias,
\emph{The $3x+1$ problem and its generalizations},
Amer. Math. Monthly \textbf{92} (1985), no.~1, 3--23.

\bibitem{lagarias-book} J.~C. Lagarias (ed.),
\emph{The Ultimate Challenge: The $3x+1$ Problem}, American Mathematical Society, 2010.

\bibitem{lothaire} M.~Lothaire,
\emph{Algebraic Combinatorics on Words}, Encyclopedia of Mathematics and its Applications
\textbf{90}, Cambridge University Press, 2002.

\bibitem{mathlib} The mathlib Community,
\emph{The Lean mathematical library}, in: Proc. 9th ACM SIGPLAN Int. Conf. on Certified
Programs and Proofs (CPP 2020), 367--381.

\bibitem{mh} M.~Morse and G.~A. Hedlund,
\emph{Symbolic dynamics II: Sturmian trajectories},
Amer. J. Math. \textbf{62} (1940), 1--42.

\bibitem{moura} L.~de Moura and S.~Ullrich,
\emph{The Lean 4 theorem prover and programming language},
in: Automated Deduction --- CADE 28, LNCS \textbf{12699}, Springer, 2021, 625--635.

\bibitem{tao} T.~Tao,
\emph{Almost all orbits of the Collatz map attain almost bounded values},
Forum Math. Pi \textbf{10} (2022), e12.

\bibitem{terras} R.~Terras,
\emph{A stopping time problem on the positive integers},
Acta Arith. \textbf{30} (1976), no.~3, 241--252.
\bibitem{bohmer} P.~E. B\"ohmer, \emph{\"Uber die Transzendenz gewisser dyadischer Br\"uche},
Math. Ann. \textbf{96} (1927), 367--377. doi:10.1007/BF01209172.
\bibitem{danilov} L.~V. Danilov, \emph{Some classes of transcendental numbers},
Math. Notes \textbf{12} (1972), 524--527. doi:10.1007/BF01095010.
\bibitem{davison} J.~L. Davison, \emph{A series and its associated continued fraction},
Proc. Amer. Math. Soc. \textbf{63} (1977), 29--32. doi:10.1090/S0002-9939-1977-0429778-5.
\bibitem{adamsdavison} W.~W. Adams and J.~L. Davison, \emph{A remarkable class of continued
fractions}, Proc. Amer. Math. Soc. \textbf{65} (1977), 194--198.
doi:10.1090/S0002-9939-1977-0441879-4.
\end{thebibliography}

\end{document}
